% ---
% title: Awards: with or without repetition
% tags: counting, permutations
% ---

\begin{problem}
	Five different awards are given to students of a class of 30. How many
	outcomes are possible if (a) a student may receive any number of awards?
	(b) each student receives at most one award?
\end{problem}

\begin{slides}
	\small
	\exstatement
	\pause
	\sol Hand out the awards one by one.
	\begin{enumerate}
		\item[(a)] Every award can go to any of the 30 students (repetition
		      allowed, order matters because the awards differ):
		      $30^5 = 24\,300\,000$.
		\item[(b)] After each award one student fewer is eligible:
		      \[ 30\cdot29\cdot28\cdot27\cdot26 = \frac{30!}{25!} = 17\,100\,720. \]
	\end{enumerate}
\end{slides}

\begin{solution}
	\idea Hand out the awards one by one. The awards are different, so it
	matters who gets which award: order matters. The two parts differ only in
	whether a student can be chosen again.

	\step{(a)} Each of the 5 awards can go to any of the 30 students:
	$30^5 = 24\,300\,000$.

	\step{(b)} The first award has 30 possible recipients, the second 29 (the
	first winner is excluded), then 28, 27, 26:
	\[ 30\cdot29\cdot28\cdot27\cdot26 = \frac{30!}{25!} = 17\,100\,720. \]
	\begin{answer}
		(a) $30^5=24\,300\,000$ \qquad (b) $30!/25!=17\,100\,720$
	\end{answer}
	\begin{remark}
		\textbf{Common mistake:} $5^{30}$ in (a). Ask ``who chooses?'': each
		\emph{award} chooses a student (30 options, 5 times), not each student an
		award.
	\end{remark}
\end{solution}
